\chapter*{Prefaci}
\addcontentsline{toc}{chapter}{Prefaci}

Aquest llibre forma part d'\textbf{aprendes}, una col·lecció de materials
d'estudi que comparteixen una mateixa manera de treballar: cada tema es
presenta en una part de \emph{teoria}, es reprèn en una part de
\emph{pràctica} amb problemes resolts, es consolida amb \emph{exercicis
proposats} i es comprova amb \emph{\nomtests}. Qui ja conegui
altres volums de la col·lecció hi reconeixerà l'estructura; qui s'hi
acosti per primera vegada la trobarà explicada a la guia de lectura.

El volum té un propòsit precís. Dins de la sèrie \textbf{logcat}, de
fonaments de la matemàtica i lògica categòrica, és el primer de quatre:
exposa la teoria de conjunts en la seva forma \emph{estructural}, el
segon, \emph{Una introducció a la teoria de categories}
\cite{perello-categories}, construeix el llenguatge general del qual
aquesta teoria és el cas fundacional, i els dos següents, dedicats a la
lògica categòrica, es basen en tots dos. Aquí els individus primitius
són les aplicacions, els conjunts són les identitats i cada construcció
---parts, sumes, productes, exponencials, nombres naturals--- es
defineix per la manera com es relaciona amb la resta mitjançant
aplicacions, i no per la naturalesa dels seus elements. És la teoria
elemental de la categoria dels conjunts de Lawvere, exposada en un
llenguatge de primer ordre que no pressuposa la teoria de categories:
per llegir aquest llibre no cal saber què és una categoria. Qui després
llegeixi el volum de categories hi reconeixerà cada construcció d'aquí
com el cas particular, en els conjunts, d'una noció general; l'apèndix
de correspondència entre els dos volums fa explícita aquesta lectura i
remet, per a cada construcció, al lloc del volum de categories on es
generalitza.

El recorregut va des dels axiomes sobre el domini, el codomini i la
composició fins a l'aritmètica dels ordinals i dels cardinals, la
cofinalitat i la hipòtesi del continu, passant per les construccions
universals, els valors de veritat, les relacions i l'elecció, les
imatges, els exponencials, el conjunt de parts i els nombres naturals.
A cada resultat s'hi indica de quins axiomes depèn, i en particular on
intervenen l'axioma d'elecció i, als darrers capítols, l'esquema de
substitució estructural que porta la teoria fins a la força de la teoria
clàssica. El capítol final fa balanç: situa els axiomes del llibre
respecte de la teoria de Zermelo--Fraenkel i de la formulació
categòrica de Lawvere.

Aquesta edició conté el volum complet: vint capítols de teoria, dinou
dels quals es desenvolupen també en les parts de pràctica, exercicis
proposats i tests, i els apèndixs, el glossari de símbols, la
bibliografia i l'índex.

Abans del primer capítol, la \emph{Introducció} explica què vol dir
mirar els conjunts des de les aplicacions i d'on ve aquesta idea.

\medskip
\noindent\textbf{Prerequisits.} El text pressuposa la maduresa matemàtica habitual d'un primer curs universitari de matemàtiques, informàtica o filosofia de la lògica: llegir fórmules de primer ordre amb quantificadors, seguir una demostració elemental i treballar informalment amb conjunts i funcions. No es pressuposa haver estudiat teoria axiomàtica de conjunts ni teoria de categories; la comparació amb la teoria clàssica es fa en observacions separades, que es poden saltar, i la terminologia categòrica apareix només com a nom de les construccions que ja s'han demostrat.

\medskip
\noindent\textbf{Errates i suggeriments.} Són benvinguts. Es poden comunicar obrint una incidència al repositori del projecte, \llibreerrates, indicant la versió (\llibreversio) i la pàgina.

% ------------------------------------------------------------
% Espai reservat per a l'autor (descomenteu i completeu):
%
% \medskip
% \noindent\textbf{Agraïments.} ...
% ------------------------------------------------------------
